\documentclass[12pt]{article}

\usepackage[a4paper, margin=2.5cm]{geometry}
\usepackage{amsmath}
\usepackage{amssymb}
\usepackage{amsthm}
\usepackage[colorlinks=true, linkcolor=blue, citecolor=blue, urlcolor=blue]{hyperref}
\usepackage{booktabs}
\usepackage{natbib}

\newtheorem{theorem}{Theorem}
\newtheorem{proposition}[theorem]{Proposition}
\newtheorem{corollary}[theorem]{Corollary}
\newtheorem{definition}[theorem]{Definition}
\newtheorem{remark}[theorem]{Remark}
\newtheorem{lemma}[theorem]{Lemma}
\newtheorem{conjecture}[theorem]{Conjecture}
\newtheorem{observation}[theorem]{Observation}

\newcommand{\Tcasc}{\mathcal{T}_{\mathrm{casc}}}
\newcommand{\Fcal}{\mathcal{F}}
\newcommand{\Zom}{\mathbb{Z}[\omega]}

\title{Caspar--Klug Triangulation Numbers from a Face-Level Cascade Operator\\[6pt]
\large Uniqueness under Local Hexagonal-Lattice Symmetry and the $T\neq5$ Correction}
\author{%
  Lee Smart\\[4pt]
  \textit{Institute of Vibrational Field Dynamics}\\[2pt]
  \texttt{contact@vibrationalfielddynamics.org}\\[2pt]
  \texttt{@vfd\_org}%
}
\date{April 2026}

\begin{document}

\maketitle

\noindent\textbf{Status:} Preprint (not peer-reviewed)

\begin{abstract}
The Caspar--Klug (1962) classification assigns each icosahedral viral
capsid a triangulation number $T = h^2 + hk + k^2$ attained by a
hexagonal-lattice subdivision of the icosahedron's twenty faces; the
integers so realised are the positive Loeschian numbers (OEIS
A003136).
Upstream cascade documents~\cite{CascadeBio} identify the biological
substrate with the binary icosahedral group $2I$ acting on the 600-cell
and observe that the dimensions of the smallest $A_5$ irreducible
representations match the smallest Caspar--Klug $T$-values, with
$T = 2$ absent from both sequences. That partial observation has two
gaps: (i)~no single cascade operator is identified whose spectrum is
the full $T$-sequence; and (ii)~the claim extends the irrep/$T$ match
to include $T = 5$, which is not a Loeschian number.

This paper isolates the classical Eisenstein norm as the canonical
face-level quadratic form under a precise $C_6$-invariance
hypothesis, notes that its positive-valued spectrum coincides with
the Caspar--Klug $T$-numbers, and corrects the $T = 5$ attribution
in upstream~\cite{CascadeBio}. No cell-level descent is proved; that
is left explicitly open. We define $\Tcasc : \Zom \to
\mathbb{Z}_{\ge 0}$ as the Eisenstein norm; its image on
$\Zom\setminus\{0\}$ is exactly the positive Loeschian numbers
(equivalently, the Caspar--Klug $T$-values), and within the class
of positive-definite quadratic forms on $\Zom$ extending to
$C_6$-invariant forms on $\mathbb{R}^2$ it is unique up to positive
scale. We also restate the small-$T$ / $A_5$-irrep
match as a precise overlap
$\{\text{distinct }A_5\text{ irrep dimensions}\}\cap
\{\text{positive Loeschians}\le 5\} = \{1, 3, 4\}$. The $T = 5$
entry in~\cite{CascadeBio} is identified as a misattribution: the
$5$-dimensional $A_5$ irrep is a summand of the face permutation
representation, not an Eisenstein-norm value, and yields no
face-level $T = 5$ value in the present model. A numerical enumeration of $C_6$-orbits up to $T \le 100$ reproduces
the listed small-$T$ multiplicities and produces one new numerical
observation in this work: $\mu(91) = 4$, the first Loeschian value
in that enumeration range admitting four distinct $C_6$-orbits,
consistent with the factorisation pattern $91 = 7 \times 13$ where
both prime factors admit chiral representatives in the wedge
picture of Section~\ref{sec:mu}.

Whether $\Tcasc$ is the face-level residue of a cell-level 600-cell
closure functional in the sense of Paper XXXII~\cite{SmartXXXII}
(pointwise) or the density-level closure form~\cite[\S F2]{CascadeFoundations}
(field) is left as an explicit open problem with a four-step plan of
attack. Two conjectural predictions for empirical capsid data are
flagged and will be tested in a follow-up note once a VIPERdb
observed-$T$ histogram is assembled. Every step is classified as
\textbf{derived} (proved), \textbf{identified} (structural
correspondence), or \textbf{open}, so the reader can locate exactly
where the chain rests on which type of support.
\end{abstract}

%==================================================================
\section{Introduction and Scope}\label{sec:intro}
%==================================================================

The Caspar--Klug theory of icosahedral virus capsids~\cite{CasparKlug}
classifies capsid subdivisions by a triangulation number
\[
    T(h, k) = h^2 + hk + k^2, \qquad (h, k) \in \mathbb{Z}^2 \setminus \{(0, 0)\},
\]
arising from replicating a hexagonal-lattice vector $h + k\omega$
(with $\omega$ a primitive sixth root of unity) across the twenty
faces of the underlying icosahedron. The integers realised by this
form are the \emph{positive} Loeschian numbers, i.e.\ the image
of the non-trivial Eisenstein-norm form on $\Zom\setminus\{0\}$.
(The Loeschian set as an arithmetic sequence includes $0 = 0^2 +
0\cdot 0 + 0^2$, but the trivial $(h, k) = (0, 0)$ subdivision is
not an admissible capsid $T$-value, so throughout this paper we
identify Caspar--Klug $T$-values with \emph{positive} Loeschian
numbers.) Characterisation: a non-negative integer $n$ is
Loeschian iff every rational prime $p \equiv 2 \pmod 3$ appearing
in its factorisation does so with even exponent; see
e.g.~\cite[Ch.\,15]{HardyWright} or~\cite[Ch.\,4]{ConwaySloane}.
The explicit positive-Loeschian sequence (OEIS A003136~\cite{OEIS-A003136})
is
\[
    1,\; 3,\; 4,\; 7,\; 9,\; 12,\; 13,\; 16,\; 19,\; 21,\; 25,\; 27,\; 28,\; 31,\;\ldots
\]
The absence of $2$, $5$, $6$, $8$, $10$, $11$, $\ldots$ from this
sequence follows from the same characterisation (each of these
non-Loeschians is divisible by an inert prime $p \equiv 2 \pmod 3$
with odd exponent).

The Vibrational Field Dynamics (VFD) cascade programme~\cite{CascadeBio}
identifies the biological substrate with the binary icosahedral group
$2I \subset \mathrm{SU}(2)$ of order $120$ acting as the vertex set of
the 600-cell, and its quotient $I = 2I/\{\pm 1\} \cong A_5$ acting as
the face-permutation group of the icosahedron. At the level of small
Caspar--Klug triangulation numbers, \cite{CascadeBio} observes:
\begin{quote}\it
The four smallest viral capsid $T$-numbers $(1, 3, 4, 5)$ are exactly
the dimensions of the four non-trivial irreducible representations of
$A_5$. And $T = 2$ is absent from both sequences \ldots\ A_5 has no
two-dimensional irrep.
\end{quote}
The quoted statement has two mathematical defects that we address in
this paper. First, $T = 5$ is \emph{not} a Loeschian number and is
therefore not a Caspar--Klug $T$-value at all. Second, the
dimension $1$ is the \emph{trivial} irreducible representation of
$A_5$, not a non-trivial one; so the phrase ``four non-trivial
irreducible representations'' conflates the four distinct non-unit
dimensions $\{3, 3', 4, 5\}$ with the four values $\{1, 3, 4, 5\}$
listed. We audit and extend the underlying observation:
\begin{enumerate}
\item[(i)] The stated match is \emph{almost} correct but includes
$T = 5$, which is \emph{not} a Loeschian number and is therefore
not a Caspar--Klug $T$-value. The honest overlap of distinct $A_5$
irrep dimensions $\{1, 3, 4, 5\}$ (four distinct values, with $1$
being the trivial representation) with the small Loeschian numbers
$\le 5$ is $\{1, 3, 4\}$ (Section~\ref{sec:t5}).
\item[(ii)] The full Caspar--Klug $T$-sequence is the spectrum of
a single face-level quadratic form on $\Zom$, unique up to positive
scale among positive-definite $C_6$-invariant quadratic forms
(Section~\ref{sec:operator}).
\item[(iii)] The cell-level extension --- whether that operator is
the face-level residue of a 600-cell closure construction --- is
left open with a four-step plan (Section~\ref{sec:open}).
\end{enumerate}

\paragraph{Scope.} All results in this paper are \emph{face-level}:
they concern the hexagonal lattice $\Zom$ living on a single
icosahedral face, with $A_5$ acting only by replicating the face-level
data across the twenty faces. We make no claim about a 600-cell
closure operator at the cell level; such an upgrade would require
establishing a principled $2I$-invariant projection that we sketch in
Section~\ref{sec:open} but do not execute.

\paragraph{Typographic conventions.} Symbol $\omega$ denotes a
primitive \emph{sixth} root of unity, $\omega = e^{i\pi/3}$, satisfying
$\omega^2 = \omega - 1$. This choice makes the Caspar--Klug norm
formula $T(h, k) = h^2 + hk + k^2$ the literal Eisenstein norm of
$h + k\omega$. The alternative convention $\omega = e^{2\pi i/3}$ (a
primitive \emph{third} root, used in~\cite[\S B3.1]{CascadeBio})
requires a basis change $k \mapsto -k$ or an alternative generator to
reproduce the plus-sign form; both attain the same Loeschian set, so
the choice is purely notational.

%==================================================================
\section{Preliminaries}\label{sec:prelim}
%==================================================================

\subsection{The icosahedral rotation group and its face representation}

$I \cong A_5$ acts on the twenty faces of the icosahedron with
point-stabiliser $C_3$ (three-fold rotation about each face's normal
axis). The permutation character $\chi_{\mathrm{perm}}$ evaluates as
\[
\chi_{\mathrm{perm}}(e) = 20,\;
\chi_{\mathrm{perm}}(3A) = 2,\;
\chi_{\mathrm{perm}}(5A) = \chi_{\mathrm{perm}}(5B) = \chi_{\mathrm{perm}}(2A) = 0,
\]
on the five $A_5$ conjugacy classes $(e, 3A, 5A, 5B, 2A)$ of sizes
$(1, 20, 12, 12, 15)$.

\begin{lemma}[$\blacktriangleright$ Decomposition of the face representation]\label{lem:L0}
The $A_5$ permutation representation on the twenty faces decomposes as
\[
\pi_{20} = \mathbf{1} \oplus \mathbf{3} \oplus \mathbf{3}' \oplus \mathbf{4}^{\oplus 2} \oplus \mathbf{5}
\]
with multiplicities $(1, 1, 1, 2, 1)$ on the irreducibles of dimensions
$(1, 3, 3, 4, 5)$. Total dimension $1 + 3 + 3 + 8 + 5 = 20$.
\end{lemma}

\begin{proof}
We apply the Frobenius inner product
$\langle \chi_{\mathrm{perm}}, \chi_\rho\rangle
= \tfrac{1}{|A_5|}\sum_c |c|\,\chi_{\mathrm{perm}}(c)\,\overline{\chi_\rho(c)}$
(the $A_5$ characters are real-valued, so the complex conjugate
reduces to $\chi_\rho(c)$). The permutation character values in
§1.2 are $\chi_{\mathrm{perm}} = (20, 2, 0, 0, 0)$ on the class
sizes $(1, 20, 12, 12, 15)$, so for any irrep $\rho$
\[
\langle \chi_{\mathrm{perm}}, \chi_\rho\rangle
= \tfrac{1}{60}\bigl(\,1\cdot 20\cdot \chi_\rho(e)
+ 20\cdot 2\cdot \chi_\rho(3A)\,\bigr)
= \tfrac{1}{3}\bigl(\chi_\rho(e) + 2\chi_\rho(3A)\bigr),
\]
since $\chi_{\mathrm{perm}}$ vanishes on the other three classes.
With the standard $A_5$ character table (e.g.~\cite{ATLAS}), whose
3A-class values are $(\chi_{\mathbf{1}}, \chi_{\mathbf{3}},
\chi_{\mathbf{3}'}, \chi_{\mathbf{4}}, \chi_{\mathbf{5}})(3A) =
(1, 0, 0, 1, -1)$, the inner products evaluate to
$(1, 1, 1, 2, 1)$ on the irreps of dimensions $(1, 3, 3, 4, 5)$.
Summation of $(\text{dim})\cdot(\text{mult})$ yields
$1 + 3 + 3 + 8 + 5 = 20$, matching $\dim\pi_{20}$.
\end{proof}

\subsection{The cascade substrate (as we use it)}

Upstream~\cite[\S 3.1]{CascadeBio} places biology at the cell
(3-volume) level of the 600-cell: $600$ tetrahedral cells. The same
document also discusses a $15$-fibre Hopf cell-fibration with $40$
cells per fibre: \cite[\S 3.2]{CascadeBio} explains the internal
arithmetic $40 = 8 \cdot 5$ of the candidate fibre size, while
\cite[\S 2.6]{CascadeBio} gives the explicit ``structural
candidate'' / ``conjectured'' status --- the fibration is
\emph{not} proved as a partition. We mirror that status. The
face-level results below do \emph{not} depend on the conjectured
Hopf fibration; the cell-level upgrade of Section~\ref{sec:open}
would.

Solidly established (\cite[\S 2.1--2.3]{CascadeBio}): $2I$ is of
order $120$ and forms the 600-cell vertex set; $I = 2I/\{\pm 1\}
\cong A_5$ is of order $60$ and acts on antipodal vertex pairs.

%==================================================================
\section{The Face-Level Cascade Operator}\label{sec:operator}
%==================================================================

\subsection{Substrate and operator}

\begin{definition}[$\blacktriangleright$ Face-level Eisenstein substrate]\label{def:D1}
The \emph{face-level icosahedral Eisenstein substrate} is the pair
\[
\mathbb{S} := (\Zom, A_5 \curvearrowright \mathcal{F}_{20}),
\]
where $\Zom$ is the ring of Eisenstein integers (hexagonal lattice on a
single icosahedral face, basis $\{1, \omega\}$ with $\omega = e^{i\pi/3}$
a primitive sixth root of unity), and $\mathcal{F}_{20}$ is the twenty-face
set carrying the $A_5$ permutation representation $\pi_{20}$ of
Lemma~\ref{lem:L0}. We model a Caspar--Klug subdivision as a single choice of
Eisenstein vector $v = h + k\omega \in \Zom \setminus \{0\}$ replicated
coherently across all twenty faces --- the same $(h, k)$ per face, up
to the local unit-ambiguity of $\Zom$.
\end{definition}

\paragraph{Two symmetries, kept distinct.} The substrate carries two
separate symmetry groups, treated throughout as independent inputs:
\begin{itemize}
\item The icosahedral group $I \cong A_5$ acts on $\mathcal{F}_{20}$
  by permuting faces, with face-stabiliser $C_3$ rotationally (or
  $D_3$ under the extended symmetry group $I_h$ of order $120$).
  \emph{Not} $C_6$.
\item The hexagonal-lattice rotation group $C_6$ acts on $\Zom$ by
  multiplication by the six Eisenstein units. This is intrinsic to
  the lattice itself, not a face-stabiliser in the icosahedral group.
\end{itemize}

\begin{definition}[$\blacktriangleright$ Face-level cascade T-operator]\label{def:D2}
The \emph{face-level cascade T-operator} is the map
\[
\Tcasc : \Zom \to \mathbb{Z}_{\ge 0},\qquad
\Tcasc(z) := z\cdot\bar{z} = |z|^2.
\]
For $z = h + k\omega$ with $\omega$ the primitive sixth root,
$\Tcasc(z) = h^2 + hk + k^2$. $\Tcasc$ is invariant under
multiplication by the six Eisenstein units (the $C_6$ action). It is
scalar-valued; no further equivariance claim is made.
\end{definition}

\subsection{Spectrum and uniqueness}

\begin{lemma}[$\blacktriangleright$ Spectrum of $\Tcasc$ = positive Loeschian numbers]\label{lem:L1}
\[
\mathrm{spec}(\Tcasc) := \Tcasc(\Zom\setminus\{0\})
= \{\text{positive Loeschian numbers}\}
= \text{Caspar--Klug}\;T\text{-values}.
\]
\end{lemma}

\begin{proof}
By construction, $\Tcasc(h + k\omega) = h^2 + hk + k^2$. The image
of this norm form on $\Zom\setminus\{0\}$ is exactly the positive
Loeschian numbers — the positive integers $n$ for which every
rational prime $p \equiv 2\pmod 3$ in the factorisation of $n$
appears with even exponent. This characterisation follows from the
structure of primes in $\Zom$, which is a Euclidean (hence
principal) domain: rational primes $p \equiv 1\pmod 3$ split as
$p = \pi\bar{\pi}$ with $N(\pi) = N(\bar{\pi}) = p$ (each a
permissible norm); rational primes $p \equiv 2\pmod 3$ are inert
with $N(\pi) = p^2$ (so any norm-factor $p$ in $n$ must appear with
even exponent); the ramified prime $3 = -\omega^2(1 - \omega)^2$
has $N(1 - \omega) = 3$. See Hardy and Wright,
\emph{An Introduction to the Theory of Numbers}~\cite{HardyWright},
Chapter~15 \S15.5 ``Primes in $\mathbb{Q}(\rho)$'', Theorems~254
and~255, for the $\Zom$ prime structure.

From this, both directions of the characterisation follow.
(\emph{Necessity.}) If $n = N(z)$ for $z \in \Zom\setminus\{0\}$,
factor $z$ into $\Zom$-primes; multiplicativity of the norm then
forces every inert-prime factor $p \equiv 2\pmod 3$ to appear with
even exponent in $n$, because each contributes $N(\pi) = p^2$.
(\emph{Sufficiency.}) Conversely, if every rational prime
$p \equiv 2\pmod 3$ in the factorisation of a positive integer $n$
appears with even exponent, then $n$ is a norm: multiply the norms
of split primes $p \equiv 1\pmod 3$ (each contributing $p = N(\pi)$
for some $\pi$ above $p$), even powers of inert primes (each
contributing $p^{2k} = N(p^k)$), and powers of the ramified prime
($3^k = N((1-\omega)^k)$), and obtain a $\Zom$-element with norm
$n$. Equivalently, see Conway and Sloane~\cite[Ch.\,4 \S4.5]{ConwaySloane}
for the Eisenstein lattice $A_2$ viewpoint. The positive-valued
enumerated sequence is OEIS A003136~\cite{OEIS-A003136}.
\end{proof}

\begin{lemma}[$\blacktriangleright$ Uniqueness under $C_6$ symmetry]\label{lem:L3}
Let $Q : \Zom \to \mathbb{R}_{\ge 0}$ be any positive-definite quadratic
form on $\Zom$ extending to a $C_6$-invariant quadratic form on
$\mathbb{R}^2 \cong \mathbb{C}$, where $C_6$ acts by rotations through
angle $\pi/3$ (equivalently, by multiplication by the six Eisenstein
units). Then $Q = c \cdot \Tcasc$ for some $c > 0$.
\end{lemma}

\begin{proof}
Write $Q(x) = x^\top M x$ with $M$ a symmetric positive-definite
$2\times 2$ matrix. Invariance under the generator $R_{\pi/3}$
(rotation by $\pi/3$) requires $R_{\pi/3}^\top M R_{\pi/3} = M$.
Writing $M = \begin{pmatrix} a & b \\ b & c\end{pmatrix}$ and
$R_{\pi/3} = \tfrac{1}{2}\begin{pmatrix} 1 & -\sqrt{3} \\ \sqrt{3} & 1 \end{pmatrix}$,
entry-by-entry comparison yields $a = c$ and $b = 0$. Equivalently,
viewing $M$ as a complex-linear operator on $\mathbb{C}$, $R_{\pi/3}$
has eigenvalues $e^{\pm i\pi/3}$ which are non-real; so the
centraliser of $R_{\pi/3}$ in the space of symmetric forms on
$\mathbb{R}^2$ is one-dimensional, consisting of scalar multiples of
the identity quadratic form. Hence $M = a\,I_2$ and $Q(x) = a|x|^2$.
Restricting to $z \in \Zom$ gives $Q(z) = a \cdot (z\bar{z}) =
a\cdot \Tcasc(z)$. Setting $c := a > 0$ completes the proof.
\end{proof}

\begin{theorem}[$\blacktriangleright$ Face-level $T$-number sequence from a unique quadratic form]\label{thm:T1}
There exists a face-level quadratic form $\Tcasc$ on $\Zom$
(Definitions~\ref{def:D1}--\ref{def:D2}), unique up to positive scale
within the class of positive-definite quadratic forms on $\Zom$ that
extend to $C_6$-invariant quadratic forms on $\mathbb{R}^2$
(Lemma~\ref{lem:L3}), whose spectrum is exactly the Caspar--Klug
$T$-number sequence (Lemma~\ref{lem:L1}).
\end{theorem}

\begin{proof}
Take $\Tcasc$ as in Definition~\ref{def:D2}. Apply Lemma~\ref{lem:L1}
for the spectrum identity and Lemma~\ref{lem:L3} for the
uniqueness-up-to-scale claim.
\end{proof}

\paragraph{Scope of Theorem~\ref{thm:T1}.} This is a face-level
statement. Whether $\Tcasc$ is the face-level residue of a
\emph{cell-level} 600-cell closure functional is the content of the
open problem in Section~\ref{sec:open}.

\paragraph{Relation to upstream~\cite[\S B3.4]{CascadeBio}.} The last
row of the B3.4 status table reads ``Full derivation of all $T$-numbers
from one cascade operator'' with status ``$\times$ open''.
Theorem~\ref{thm:T1} isolates the unique face-level quadratic
form, among those extending to $C_6$-invariant forms on
$\mathbb{R}^2$, with the Caspar--Klug spectrum. This is a
normal-form / canonical-operator statement at face level and does
\emph{not} constitute a descent from the cell-level 600-cell
machinery; the latter remains open per Section~\ref{sec:open}.

%==================================================================
\section{Settling the $T = 5$ Question}\label{sec:t5}
%==================================================================

\begin{remark}[$\triangleright$ The upstream $T = 5$ entry is misleading]\label{rem:R1}
The B3.2 table in~\cite[\S B3.2]{CascadeBio} lists
\[
\text{``}T = 5 \;\leftrightarrow\; A_5\ \text{5-dim irrep''}
\]
as part of the small-$T$ / $A_5$-irrep match. This entry is not
supported by either of the following two mathematical facts:
\begin{enumerate}
\item[(i)] The classical Caspar--Klug sequence: $5$ is not a Loeschian
number (no $(h, k) \in \mathbb{Z}^2$ satisfies $h^2 + hk + k^2 = 5$,
because $5$ is inert in $\Zom$ with odd exponent in its factorisation,
equivalently $5 \equiv 2 \pmod 3$).
\item[(ii)] The spectrum of $\Tcasc$: by Lemma~\ref{lem:L1}, $5
\notin \mathrm{spec}(\Tcasc)$.
\end{enumerate}
What the $A_5$ 5-dim irrep \emph{is} is a summand of the face
permutation rep $\pi_{20}$ (Lemma~\ref{lem:L0}); it is a
representation-theoretic invariant, not an Eisenstein-norm value,
and yields no face-level $T = 5$ value in the present model. (The
empirical question of whether observed capsid counts additionally
reinforce this absence is distinct and is deferred to a follow-up
test against a VIPERdb observed-$T$ histogram.)
\end{remark}

\begin{lemma}[$\blacktriangleright$ Honest small-$T$ / $A_5$ match]\label{lem:L2}
$A_5$ has five irreducible representations with \emph{four} distinct
dimensions: $1, 3, 4, 5$; the dimension $3$ occurs twice, for a
Galois-conjugate pair usually written $\mathbf{3}$ and $\mathbf{3}'$
(swapped under $\varphi \leftrightarrow 1 - \varphi$). The Loeschian
numbers $\le 5$ are $\{1, 3, 4\}$. Hence
\[
\{\text{distinct } A_5 \text{ irrep dims}\} \cap \{\text{Loeschians} \le 5\} = \{1, 3, 4\}.
\]
Three of the four distinct $A_5$ irrep dimensions coincide with small
Loeschian numbers. The dimension $5$ is not Loeschian; the duplicate
dimension in $\mathbf{3}'$ carries the same value as $\mathbf{3}$ and
contributes no new $T$-value.
\end{lemma}

\begin{proof}
First claim: Lemma~\ref{lem:L0} and the standard $A_5$ character
table (e.g.~\cite{ATLAS}). Second claim: among positive integers
$n\le 5$, the Loeschian criterion selects exactly those $n$ in
whose prime factorisation every prime $p\equiv 2\pmod 3$ appears
with even exponent. This excludes $2$ (prime, inert, odd exponent)
and $5$ (prime, inert, odd exponent), leaving $\{1, 3, 4\}$.
Intersecting with the distinct $A_5$ irrep dimensions $\{1, 3, 4, 5\}$
yields $\{1, 3, 4\}$.
\end{proof}

\paragraph{Downstream edit.} \cite[\S B3.2]{CascadeBio}'s table should
be amended per Remark~\ref{rem:R1} and Lemma~\ref{lem:L2}. We do not
rewrite the upstream document here; the edit is flagged for a future
revision of the bio document.

\paragraph{Galois remark.} $\mathbf{3}'$ has the same dimension as
$\mathbf{3}$, contributing no new $T$-value, but character values on
the 5-cycle classes are Galois-conjugate swapped under
$\varphi \leftrightarrow 1 - \varphi$. Conjecture~\ref{conj:C4}
in~\S\ref{sec:open} records the possibility of a corresponding
structural dimorphism at $T = 3$; no claim is made here.

%==================================================================
\section{Orbit Multiplicity and Chirality}\label{sec:mu}
%==================================================================

\begin{definition}[$\blacktriangleright$ Orbit multiplicity $\mu(T)$]\label{def:D3}
For each positive Loeschian integer $T$,
\[
\mu(T) := \#\{\, C_6\text{-orbits of } (h, k) \in \mathbb{Z}^2\setminus\{0\}
\text{ with } h^2 + hk + k^2 = T\,\},
\]
where $C_6$ acts on $\Zom$ by multiplication by the six Eisenstein
units. We do \emph{not} quotient by the reflection $h \leftrightarrow
k$: this is a modelling choice, treating mirror-related wedge
representatives as distinct oriented face-level classes, so that the
resulting orbit structure carries the laevo/dextro (left-handed /
right-handed) distinction present in empirical capsid catalogues.
(The ambient icosahedral rotation group $I \cong A_5$ is a proper
rotation group containing no reflections, which makes this
convention structurally consistent with the face-level substrate
but does not by itself prove that it is the correct empirical
choice; that is a modelling assumption tested downstream by
Conjecture~\ref{conj:C3}.)
\end{definition}

\paragraph{Sample values} (verified by direct enumeration; see
Section~\ref{sec:num}):
\begin{align*}
\mu(1) &= \mu(3) = \mu(4) = \mu(9) = \mu(12) = \mu(16) = \mu(25) = \mu(27) = 1,\\
\mu(7) &= \mu(13) = \mu(19) = \mu(21) = \mu(28) = \mu(31) = \mu(37) = \mu(39) = 2,\\
\mu(49) &= 3 \quad\text{representatives } (7,0), (5,3), (3,5),\\
\mu(91) &= 4 \quad\text{representatives } (1,9), (5,6), (6,5), (9,1).
\end{align*}

\begin{observation}[Heuristic structure; direct enumeration verifies the listed finite cases]\label{obs:mu-structure}
$\mu(T)$ counts distinct $C_6$-orbits at $T$. Heuristically, a single
norm-factorisation contributes one orbit ($\mu = 1$) when its
canonical wedge representative lies on a hexagonal-lattice reflection
axis --- either $k = 0$ or $h = k$ --- because such an orbit coincides
with its mirror image under the $h \leftrightarrow k$ reflection
(achiral, laevo and dextro indistinguishable). A factorisation with
a wedge representative in the open positive sextant ($h, k > 0$,
$h \neq k$) contributes two orbits ($\mu = 2$), a laevo/dextro
chiral pair. $T$ with multiple independent norm-factorisations can
yield $\mu(T) \ge 3$: $T = 49 = 7^2$ gives $\mu = 3$ (one axis
orbit from $(7, 0)$ plus one chiral pair from $7\times 7$);
$T = 91 = 7\cdot 13$ gives $\mu = 4$ (two chiral pairs, one per prime
factor's contribution). This is heuristic explanation, not a proved
classification; the direct enumeration in Section~\ref{sec:num}
verifies the listed finite cases.
\end{observation}

\begin{conjecture}[$\circ$ Closure-preference weighting]\label{conj:C1}
Let $Z(T)$ denote a cascade closure-preference weight on capsid
$T$-numbers at equilibrium. Then
\[
Z(T) \propto \mu(T) \cdot \exp\!\big(-\beta \cdot c(T)\big),
\]
where $c(T)$ is a monotone non-decreasing assembly cost and $\beta$
a positive cascade constant of order unity. The $\mu(T)$ factor is
the natural degeneracy weight once one adopts an equal-weight-per-
$C_6$-orbit modelling assumption; it is not forced by
Definition~\ref{def:D3} itself, since a definition of an orbit-count
does not determine a probability weight without an additional
weighting choice. The exponential ansatz on $c(T)$ is a separate
modelling choice. Status: untested; test deferred to a follow-up
comparison against a VIPERdb observed-$T$ histogram.
\end{conjecture}

\begin{conjecture}[$\circ$ Chirality asymmetry from $2I$ handedness]\label{conj:C3}
For chiral $T$-numbers (those with $\mu(T) \ge 2$ and a
non-axis-aligned representative), the cascade predicts a non-zero
laevo/dextro empirical asymmetry in observed capsid counts,
\[
\frac{N_{\mathrm{laevo}} - N_{\mathrm{dextro}}}{N_{\mathrm{laevo}} + N_{\mathrm{dextro}}} \neq 0,
\]
with sign set by the ambient biological chirality (L-amino-acid
world). Magnitude not derived from first principles. Extends the $2I$
handedness / L-amino-acid-preference thread in~\cite[\S 2.7]{CascadeBio}
from the amino-acid scale to the capsid scale. Qualitative only.
\end{conjecture}

%==================================================================
\section{Open Problem: Relation to the 600-Cell Closure Machinery}\label{sec:open}
%==================================================================

We ask whether $\Tcasc$ is the face-level residue of a cell-level
600-cell closure construction. The question splits by which upstream
closure object is the target of the restriction.

\paragraph{The two candidate upstream objects.}
\begin{itemize}
\item \textbf{Paper XXXII~\cite{SmartXXXII}, \S 4, Theorem
``Uniqueness of the closure functional'':} within the class
$\mathcal{C}$ of continuous regularised soft-min functionals
$\Fcal^V_\varepsilon(x) = \mathrm{opt}_p\!\left[
\sum_i p_i\, |x - v_i|^2/2 + \varepsilon^2 R(p)\right]$ with
permutation-symmetric $R$ satisfying $R(\delta_i) = 0$ at each
vertex simplex and $\mathrm{opt} \in \{\min, \max\}$ (all part of
the class definition), axioms~\textbf{A1} continuity, \textbf{A2}
sharp limit, and \textbf{A3} extensivity under independent
composition force, up to reparametrisation of $\varepsilon$ and
additive constants,
$\Fcal_\varepsilon(x) = -\varepsilon^2 \log \sum_i \exp(-|x - v_i|^2/(2\varepsilon^2))$,
i.e.~the log-sum-exp closure functional with Gibbs minimiser. $\Fcal$
is \emph{pointwise} on the manifold: $\Fcal^V_\varepsilon : M \to
\mathbb{R}$.
\item \textbf{\cite[\S F2]{CascadeFoundations}}: density-level
closure functional $\Fcal[\Phi] = \int (\alpha R + \beta E - \gamma Q)\,dV$,
the unique form \emph{up to the coefficients} $\alpha, \beta, \gamma$
under locality, $W(E_8)$-/Galois-twist-/translation invariance,
order at most $2$, and scalar-valuedness. This is a \emph{field}
functional on $\Phi \in C^\infty(\cdot, \mathcal{A})$ for a
suitable target bundle $\mathcal{A}$.
\end{itemize}
The two upstream objects are related but not identical; an earlier
draft of this work mistakenly attributed field-theoretic properties
to Paper XXXII's pointwise $\Fcal$. We now keep them distinct.

\begin{conjecture}[$\circ$ Face-level residue; open problem]\label{conj:C2}
The conjecture splits into two interpretations:
\begin{itemize}
\item[(C.2-a)] \emph{Point-wise form}: is there a principled
$2I$-invariant restriction of $\Fcal^V_\varepsilon : S^3 \to \mathbb{R}$
to a face-level operator on $\Zom$? As stated this is ill-posed --- a
pointwise function on $S^3$ does not directly give a quadratic form
on the face lattice --- so (C.2-a) reduces in practice to (C.2-b)
after a field-theoretic lift.
\item[(C.2-b)] \emph{Field / density form}: does the density-level
closure functional $\Fcal[\Phi] = \int(\alpha R + \beta E - \gamma Q)\,dV$
of~\cite[\S F2]{CascadeFoundations}, restricted to $2I$-invariant
fields and face-level-averaged via a principled Hopf cell-fibre map,
reduce on each face to $\lambda\cdot\Tcasc$ for some $\lambda > 0$?
\end{itemize}
A four-step plan of attack for (C.2-b):
\begin{enumerate}
\item Identify the $2I$-invariant subspace of F2 field configurations
on which $\Fcal[\Phi]$ restricts coherently.
\item Construct a principled face-level projection from that subspace
to $\Zom$-valued face data. A natural candidate average is the
conjectured Hopf cell-fibre partition, with fibre arithmetic
$40 = 8 \cdot 5$ per~\cite[\S 3.2]{CascadeBio} and the explicit
conjectural / structural-candidate status per~\cite[\S 2.6]{CascadeBio};
any face-level projection built on it inherits conjectural status
until the fibration is established. Alternative averaging schemes
may exist.
\item Show $\Fcal[\Phi]$ restricts along this projection to a
positive-definite quadratic form on $\Zom$. F2 gives a quadratic
density, but quadratic density does not automatically yield
quadratic-on-slice after averaging; this step is the real technical
obstacle.
\item Apply Lemma~\ref{lem:L3} to identify the projected form with
$c \cdot \Tcasc$ for some $c > 0$.
\end{enumerate}
Only step~(4) is established here; steps~(1)--(3) are open. The
conjecture is therefore recorded as an open problem with a concrete
plan, not as a proved descent theorem.
\end{conjecture}

\begin{conjecture}[$\circ$ Galois-conjugate structural dimorphism at $T = 3$]\label{conj:C4}
The Galois pair $\mathbf{3}, \mathbf{3}'$ of $A_5$ (Lemma~\ref{lem:L0})
may manifest as a discrete structural feature of $T = 3$ capsids:
coat-protein orientation dimorphism, golden-ratio / silver-ratio
radius bistability, or a pair of near-degenerate conformations
separated by the Galois symmetry. Speculative; no experimental handle
proposed. Recorded for future attention.
\end{conjecture}

%==================================================================
\section{Numerical Verification and Planned Empirical Test}\label{sec:num}
%==================================================================

\subsection{Enumeration output (present)}

A Python enumeration script at
\texttt{papers/biology-rung/scripts/wo1\_capsid\_predict.py}
enumerates $(h, k) \in \mathbb{Z}^2$ within the non-negative wedge,
generates $C_6$-orbits via the action $(h, k)\mapsto(-k, h+k)$, and
emits a per-orbit CSV at
\texttt{papers/biology-rung/data/wo1\_prediction.csv}. With default
cutoff $T \le 100$:

\begin{itemize}
\item $36$ distinct Loeschian $T$-values $\le 100$;
\item $61$ distinct $C_6$-orbits summed over those $T$;
\item Direct enumeration up to $T \le 100$ gives $\mu(91) = 4$,
with orbit representatives $(1, 9), (5, 6), (6, 5), (9, 1)$ --- the
first Loeschian value up to $T \le 100$ with four $C_6$-orbits.
This is consistent with the factorisation pattern $91 = 7 \cdot 13$
where both prime factors admit chiral representatives in the wedge
picture (heuristic reading, not a proved classification theorem).
\end{itemize}

All spot-check assertions against the sample $\mu(T)$ values
following Definition~\ref{def:D3} and Observation~\ref{obs:mu-structure}
pass ($18$ checks); all forbidden-$T$ absence checks pass
($11$ checks on $T \in \{2, 5, 6, 8, 10, 11, 14, 15, 17, 18, 20\}$).

\subsection{Planned empirical test (VIPERdb, future work)}

A companion script \texttt{wo1\_viperdb\_fetch.py} provides two data
modes, \texttt{-{}-fetch} (public VIPERdb API) and
\texttt{-{}-from-literature} (curated CSV). At the time of writing,
the four candidate VIPERdb endpoints coded into the fetcher return
HTTP 404 (VIPERdb's API has evolved past the attempted URLs), so an
empirical observed-$T$ histogram has not yet been assembled. The
comparison script \texttt{wo1\_compare.py} is itself present and
operational; when invoked with no observed dataset it emits a
\texttt{WAITING\_FOR\_DATA} report at
\texttt{data/wo1\_comparison.md} rather than fabricating output. A
follow-up note, once the observed dataset is obtained (endpoint
refresh or literature-mode population), will carry out the
$\chi^2$ / KL / chirality-asymmetry tests against three successive
null hypotheses:
\begin{itemize}
\item[$H_0$] $T$ uniform over positive integers (expected strongly
falsified by the Loeschian mathematical confinement);
\item[$H_1$] $T$ uniform over Loeschians (expected weakly falsified
by empirical over-representation of small $T$);
\item[$H_2$] $T$ distributed as
$Z(T) \propto \mu(T)\exp(-\beta c(T))$ per Conjecture~\ref{conj:C1},
with $c(T)$ either a monotone ansatz or an empirical fit.
\end{itemize}
No empirical claim is made in the present paper; this section
describes the test design, not the result.

%==================================================================
\section{Summary and Logical Chain}\label{sec:chain}
%==================================================================

\paragraph{What is derived ($\blacktriangleright$).}
\begin{itemize}
\item Theorem~\ref{thm:T1}: face-level quadratic form $\Tcasc$ on
$\Zom$ with $\mathrm{spec}(\Tcasc) = $ Loeschian numbers, unique up
to positive scale within the class of positive-definite quadratic
forms on $\Zom$ extending to $C_6$-invariant forms on $\mathbb{R}^2$.
\item Lemma~\ref{lem:L3}: matrix-centraliser uniqueness proof, via
the matrix centraliser of the $C_6$ rotation.
\item Lemma~\ref{lem:L1}: spectrum identity, from the classical
Eisenstein-norm characterisation of Loeschian integers.
\item Lemma~\ref{lem:L0}: $\pi_{20}$ decomposition via Frobenius
inner product on $A_5$.
\item Remark~\ref{rem:R1} and Lemma~\ref{lem:L2}: $T = 5$ entry
in~\cite[\S B3.2]{CascadeBio} is mathematically unsupported;
the honest overlap is $\{1, 3, 4\}$.
\end{itemize}

\paragraph{What is identified ($\triangleright$).}
\begin{itemize}
\item The face-level $C_6$ is the hexagonal lattice's own rotational
symmetry, not an icosahedral face stabiliser (which is only $C_3$).
\item $\Tcasc$ is the classical Eisenstein norm; the cascade-specific
content is its uniqueness under $C_6$ (Lemma~\ref{lem:L3}) and its
identification as the face-level candidate for a cell-level cascade
operator.
\item New numerical observation in this work: $\mu(91) = 4$, the
first Loeschian $T$ up to $T \le 100$ with four distinct
$C_6$-orbits ($T = 91$ itself; exhaustive enumeration over that
range).
\end{itemize}

\paragraph{What is open ($\circ$).}
\begin{itemize}
\item Conjecture~\ref{conj:C2}: whether $\Tcasc$ descends from a
cell-level 600-cell closure functional (Paper XXXII~\cite{SmartXXXII}
or~\cite[\S F2]{CascadeFoundations}). Four-step plan; only step~(4)
established.
\item Conjecture~\ref{conj:C1}: closure-preference weighting ansatz
and the form of the assembly cost $c(T)$. Untested.
\item Conjecture~\ref{conj:C3}: magnitude of the chirality asymmetry.
Sign predicted qualitatively, magnitude not derived.
\item Conjecture~\ref{conj:C4}: Galois-conjugate structural
dimorphism at $T = 3$. Speculative.
\item Empirical: $\chi^2$ / KL / chirality tests against a VIPERdb
observed-$T$ histogram. The comparison script
\texttt{wo1\_compare.py} is present and operational; the remaining
blocker is the observed-$T$ dataset itself (candidate VIPERdb
endpoints returned 404 at the time of writing; see §7).
\item Downstream edit: \cite[\S B3.2]{CascadeBio} table revision per
Remark~\ref{rem:R1}.
\end{itemize}

%==================================================================
\section{Reproducibility}\label{sec:repro}
%==================================================================

This manuscript lives at
\texttt{papers/biology-rung-capsid/biology-rung-capsid.tex}; the
accompanying mathematical and computational artefacts (derivation
notes, math-catalogue ledger, enumeration and comparison scripts,
data artefacts, and the WO spec) are all kept in the sibling
directory \texttt{papers/biology-rung/} of the VFD repository to
reflect the umbrella-programme structure. All claims and numerical
statements in this paper can be reproduced from the artefacts in
\texttt{papers/biology-rung/}:
\begin{itemize}
\item \texttt{derivation-capsid.md}: Markdown derivation source from
which this paper is distilled.
\item \texttt{math-catalogue.md}: ledger of every $D_i, L_i, T_i,
C_i, R_i, N_i$ object introduced, with provenance and status.
\item \texttt{scripts/wo1\_capsid\_predict.py}: orbit enumeration
and spot-check script. Run with \texttt{python3 wo1\_capsid\_predict.py}
(default cutoff $T \le 100$).
\item \texttt{data/wo1\_prediction.csv}: canonical per-orbit output
with columns \texttt{T, mu, representative, chirality}.
\item \texttt{scripts/wo1\_viperdb\_fetch.py}: fetch infrastructure
(API mode and curated-CSV mode) for the future empirical comparison.
\item \texttt{WO-BIOLOGY-RUNG-001-CAPSID.md}: work order specifying
acceptance criteria and open items.
\end{itemize}

\begin{thebibliography}{9}

\bibitem{CasparKlug}
Caspar, D.~L.~D.\ and Klug, A.,
\emph{Physical Principles in the Construction of Regular Viruses},
Cold Spring Harbor Symposia on Quantitative Biology \textbf{27}
(1962), 1--24.

\bibitem{OEIS-A003136}
OEIS Foundation Inc.,
\emph{The On-Line Encyclopedia of Integer Sequences}, Sequence
A003136 (Loeschian numbers), \url{https://oeis.org/A003136}.

\bibitem{HardyWright}
Hardy, G.~H.\ and Wright, E.~M.,
\emph{An Introduction to the Theory of Numbers}, 6th ed.\ (revised
by D.~R.~Heath-Brown and J.~H.~Silverman), Oxford University
Press, 2008. Chapter~15 (quadratic fields) contains the
characterisation of integers representable as $h^2 + hk + k^2$ via
the arithmetic of $\mathbb{Z}[\omega]$.

\bibitem{ConwaySloane}
Conway, J.~H.\ and Sloane, N.~J.~A.,
\emph{Sphere Packings, Lattices and Groups}, 3rd ed., Springer
Grundlehren der mathematischen Wissenschaften \textbf{290}, 1999.
Chapter~4 (§4.5) and Chapter~7 cover the Eisenstein lattice $A_2$
and its norm form.

\bibitem{ATLAS}
Conway, J.~H., Curtis, R.~T., Norton, S.~P., Parker, R.~A., and
Wilson, R.~A.,
\emph{ATLAS of Finite Groups}, Oxford University Press, 1985.

\bibitem{CascadeBio}
\emph{Cascade Biology Rung: The 600-Cell Cell-Level Substrate},
internal manuscript,
\texttt{papers/cascade-derivation/cascade-bio.md}. Preprint, Institute
of Vibrational Field Dynamics.

\bibitem{CascadeFoundations}
\emph{Cascade Foundations: F1--F8}, internal manuscript,
\texttt{papers/cascade-derivation/cascade-foundations.md}. Preprint,
Institute of Vibrational Field Dynamics.

\bibitem{SmartXXXII}
Smart, L.,
\emph{The Closure Functional from First Principles: Self-Similar
Permeability, the 600-Cell, and the Derived Interaction},
\texttt{papers/paper-xxxii/paper-xxxii.tex} (Paper XXXII).
Preprint, Institute of Vibrational Field Dynamics, 2026.

\end{thebibliography}

\end{document}
